% ID: SEL-D0001-OB-E01
% ESTADO: validada. S002, 2026-11-07. Tres pasadas del mismo autor; ver informe.
\textbf{a) Simetría.}\par\nopagebreak[4]

Una matriz cuadrada es simétrica si coincide con su traspuesta: $a_{ij}=a_{ji}$ para todos los pares de índices. En este caso
\[
a_{23}=-1\ne1=a_{32},
\]
por lo que \textbf{$A$ no es simétrica}. Basta un par de entradas distintas para descartar la simetría.

\textbf{b) Cálculo de la inversa.}\par\nopagebreak[4]

Una matriz cuadrada tiene inversa si su determinante es distinto de cero. Desarrollando por la primera fila,
\[
\det A=1\begin{vmatrix}2&-1\\1&1\end{vmatrix}
-(-2)\begin{vmatrix}-2&-1\\0&1\end{vmatrix}
=3-4=-1.
\]
El cofactor $C_{ij}$ se obtiene eliminando la fila $i$ y la columna $j$, calculando el determinante restante y aplicando el signo $(-1)^{i+j}$. La matriz de cofactores y su traspuesta son
\[
C=\begin{pmatrix}3&2&-2\\2&1&-1\\2&1&-2\end{pmatrix},
\qquad
\operatorname{adj}(A)=C^t=\begin{pmatrix}3&2&2\\2&1&1\\-2&-1&-2\end{pmatrix}.
\]
Así,
\[
\boxed{A^{-1}=\frac{\operatorname{adj}(A)}{\det A}
=\begin{pmatrix}-3&-2&-2\\-2&-1&-1\\2&1&2\end{pmatrix}.}
\]
Se comprueba multiplicando: $AA^{-1}=A^{-1}A=I_3$.

\textbf{c) Ecuación matricial.}\par\nopagebreak[4]

En el enunciado $O$ representa la matriz nula de orden 3. Aislamos el producto que contiene la incógnita:
\[
2XA=A^2+3I_3.
\]
Como $A$ multiplica a $X$ por la derecha, multiplicamos ambos miembros por $A^{-1}$ también por la derecha. El orden de los factores es esencial:
\[
2XAA^{-1}=(A^2+3I_3)A^{-1}
\quad\Longrightarrow\quad
X=\frac12(A+3A^{-1}).
\]
La simplificación usa asociatividad y $AA^{-1}=I_3$, sin suponer que matrices arbitrarias conmuten. Sustituyendo los valores,
\[
A+3A^{-1}=\begin{pmatrix}-8&-8&-6\\-8&-1&-4\\6&4&7\end{pmatrix},
\qquad
\boxed{X=\begin{pmatrix}-4&-4&-3\\-4&-\frac12&-2\\3&2&\frac72\end{pmatrix}.}
\]
\textbf{Comprobación.} El cuadrado de $A$ y el producto obtenido son
\[
A^2=\begin{pmatrix}5&-6&2\\-6&7&-3\\-2&3&0\end{pmatrix},\qquad
2XA=\begin{pmatrix}8&-6&2\\-6&10&-3\\-2&3&3\end{pmatrix}=A^2+3I_3.
\]
Por tanto, $2XA-A^2-3I_3=O$, exactamente como exige la ecuación.
